\section{The Appointed Texts}
\label{sec:appointed-texts}

The Latin below is that of the 1962 typical edition, \work{Missale Romanum}
(Typis Polyglottis Vaticanis), printed pp.~702--703, nos.~3974--3983, for the
feast, and pp.~416--417, nos.~1720, 1726 and 1728, for the three orations of
the Twenty-third Sunday after Pentecost, read on the page images of the CMAA
facsimile. Every word of the feast's ten texts stands in the Tours printing of
the \work{Missale Romanum} (Mame, 1922), pp.~748--750, and every word of the
Sunday's three orations in the Ratisbon printing (Pustet, 1862),
pp.~357--358. Those printings differ from the 1962 page in orthography (the
Mame writes \latin{Joannes}, \latin{Jesus}, \latin{justi}), in pointing, and
in the length of some conclusion cues; no word that is said differs. The
rubrics of the 1962 book are described here, not reproduced.

The English of the scriptural texts is the Douay--Rheims Bible in Bishop
Challoner's revision, which translates the Clementine Vulgate in the Missal's
own numbering. It is given at the canonical verses: where a chant takes part
of a verse, the Douay words for that part are given and an ellipsis marks what
the Missal leaves out, and each place where the Missal's form departs from the
Vulgate is named beneath the text. The English of the Introit antiphon and of
the six orations is that of the Lasance \work{New Roman Missal} in the
Benziger revision of 1945, pp.~1265 and 1268 for the feast and pp.~733 and
735--736 for the Sunday. The Lasance prints this Mass under the rubrics in
force before 1960, as a double of the first class with commemorations
arranged differently; only its English is used here, and none of this English
is an approved liturgical translation of the 1962 Missal. Psalms are cited in
the Missal's numbering, with the Hebrew numbering in parentheses.

\subsection{Introit}
\label{proper-introit}

\properlocus{\latin{Antiphona ad Introitum}; verse \latin{Ps.\ 32, 1} (Heb.\ 33). Missal p.~702, no.~3974.}

\begin{latinproper}
Gaudeámus omnes in Dómino, diem festum celebrántes sub honóre Sanctórum
ómnium: de quorum solemnitáte gaudent Angeli, et colláudant Fílium Dei.
\textnormal{Ps.\ 32, 1} Exsultáte, iusti, in Dómino: rectos decet collaudátio.
℣.~Glória Patri.
\end{latinproper}

\begin{englishwitness}{Lasance, The New Roman Missal (1945), All Saints, Introit}
Let us all rejoice in the Lord, celebrating a feast in honor of all the
saints, in whose solemnity the angels rejoice and join in praising the Son of
God. Rejoice in the Lord, ye just: praise becometh the upright. V. Glory.
\end{englishwitness}

\begin{englishwitness}{Douay--Rheims (Challoner), Ps 32:1, the verse}
\dots\ Rejoice in the Lord, O ye just: praise becometh the upright.
\end{englishwitness}

\noindent The antiphon is not taken from Scripture; it is a text of the
Church, and the Missal gives it no citation. The verse is the first of
Psalm~32. The Douay prints the psalm's inscription, ``A psalm for David,''
inside its verse~1, and the Missal does not sing it. After the \latin{Gloria
Patri} the antiphon is repeated (RGMR 427).

\subsection{Collect}
\label{proper-collect}

\properlocus{\latin{Oratio}. Missal p.~702, no.~3975.}

\begin{latinproper}
Omnípotens sempitérne Deus, qui nos ómnium Sanctórum tuórum mérita sub una
tribuísti celebritáte venerári: quǽsumus; ut desiderátam nobis tuæ
propitiatiónis abundántiam, multiplicátis intercessóribus, largiáris. Per
Dóminum.
\end{latinproper}

\begin{englishwitness}{Lasance, The New Roman Missal (1945), All Saints, Collect}
Almighty, eternal God, Who hast granted us to venerate the merits of all Thy
saints in one solemnity, vouchsafe to us, we beseech Thee, through the
multitude of our intercessors, that abundance of Thy mercy for which we ever
have hoped. Through our Lord.
\end{englishwitness}

\noindent The Latin's \latin{sub una \dots\ celebritate}, ``in one
celebration,'' is the Lasance's ``in one solemnity''; its
\latin{multiplicatis intercessoribus}, an ablative absolute in which the
intercessors are multiplied, is the Lasance's ``through the multitude of our
intercessors.'' The prayer is addressed to the Father, and the Missal's cue
\latin{Per Dominum} stands for the full conclusion through the Son.

\subsection{Commemoration Collect}
\label{proper-commemoration-collect}

\properlocus{\latin{Oratio} of the Twenty-third Sunday after Pentecost. Missal p.~416, no.~1720.}

\begin{latinproper}
Absólve, quǽsumus, Dómine, tuórum delícta populórum: ut a peccatórum
néxibus, quæ pro nostra fragilitáte contráximus, tua benignitáte liberémur.
Per Dóminum nostrum Iesum Christum, Fílium tuum: Qui tecum vivit et regnat
in unitáte.
\end{latinproper}

\begin{englishwitness}{Lasance, The New Roman Missal (1945), Twenty-third Sunday after Pentecost, Collect}
Remit, we beseech Thee, O Lord, the sins of Thy people, that by Thy kindness
we may be delivered from the trammels of our sins, in which, through our
frailty, we have become entangled. Through our Lord.
\end{englishwitness}

\noindent The priest says this prayer after the feast's Collect, as a
commemoration of the Sunday, under its own conclusion (RG 108, 109~a, 111~a;
RGMR 434~a, 437). The Latin's \latin{populorum} is plural, ``of thy peoples'';
the Lasance gives the singular.

\subsection{Lesson}
\label{proper-lesson}

\properlocus{\latin{Lectio libri Apocalypsis beati Ioannis Apostoli}, Apoc 7:2--12. Missal pp.~702--703, no.~3976.}

\begin{latinproper}
In diébus illis: Ecce ego Ioánnes vidi álterum Angelum ascendéntem ab ortu
solis, habéntem signum Dei vivi: et clamávit voce magna quátuor Angelis,
quibus datum est nocére terræ et mari, dicens: Nolíte nocére terræ et mari,
neque arbóribus, quoadúsque signémus servos Dei nostri in fróntibus eórum.
Et audívi númerum signatórum, centum quadragínta quátuor mília signáti, ex
omni tribu filiórum Israel. Ex tribu Iuda duódecim mília signáti. Ex tribu
Ruben duódecim mília signáti. Ex tribu Gad duódecim mília signáti. Ex tribu
Aser duódecim mília signáti. Ex tribu Néphtali duódecim mília signáti. Ex
tribu Manásse duódecim mília signáti. Ex tribu Símeon duódecim mília
signáti. Ex tribu Levi duódecim mília signáti. Ex tribu Issachar duódecim
mília signáti. Ex tribu Zábulon duódecim mília signáti. Ex tribu Ioseph
duódecim mília signáti. Ex tribu Béniamin duódecim mília signáti. Post hæc
vidi turbam magnam, quam dinumeráre nemo póterat, ex ómnibus géntibus, et
tríbubus, et pópulis, et linguis: stantes ante thronum, et in conspéctu
Agni, amícti stolis albis, et palmæ in mánibus eórum: et clamábant voce
magna, dicéntes: Salus Deo nostro, qui sedet super thronum, et Agno. Et
omnes Angeli stabant in circúitu throni, et seniórum, et quátuor
animálium: et cecidérunt in conspéctu throni in fácies suas, et
adoravérunt Deum, dicéntes: Amen. Benedíctio, et cláritas, et sapiéntia, et
gratiárum áctio, honor, et virtus, et fortitúdo Deo nostro in sǽcula
sæculórum. Amen.
\end{latinproper}

\begin{englishwitness}{Douay--Rheims (Challoner), Apoc 7:2--12}
And I saw another angel ascending from the rising of the sun, having the sign
of the living God. And he cried with a loud voice to the four angels to whom
it was given to hurt the earth and the sea, Saying: Hurt not the earth nor the
sea nor the trees, till we sign the servants of our God in their foreheads.
And I heard the number of them that were signed. An hundred forty-four
thousand were signed, of every tribe of the children of Israel. Of the tribe
of Juda, twelve thousand signed: \dots\ Ruben \dots\ Gad \dots\ Aser \dots\
Nephthali \dots\ Manasses \dots\ Simeon \dots\ Levi \dots\ Issachar \dots\
Zabulon \dots\ Joseph \dots\ Benjamin \dots\ After this, I saw a great
multitude, which no man could number, of all nations and tribes and peoples
and tongues, standing before the throne and in sight of the Lamb, clothed with
white robes, and palms in their hands. And they cried with a loud voice,
saying: Salvation to our God, who sitteth upon the throne and to the Lamb. And
all the angels stood round about the throne and the ancients and the four
living creatures. And they fell down before the throne upon their faces and
adored God, Saying: Amen. Benediction and glory and wisdom and thanksgiving,
honour and power and strength, to our God, for ever and ever. Amen.
\end{englishwitness}

\noindent For each of the twelve tribes after Juda the Douay repeats ``twelve
thousand signed''; the repetitions are shortened here to the tribe's name. The
Missal's opening, \latin{In diebus illis: Ecce ego Ioannes vidi alterum
Angelum}, stands for the Vulgate's \latin{Et vidi alterum angelum}, the
Douay's ``And I saw another angel.'' The Lasance renders the Missal's opening
``In those days, behold, I, John, saw another angel ascending from the rising
of the sun.'' Every other word of verses 2--12 agrees with the Clementine
text. The Lesson stops
before verse~13, where one of the ancients asks John who the white-robed
multitude are and answers that they ``are come out of great tribulation''
(7:13--14).

\subsection{Gradual}
\label{proper-gradual}

\properlocus{\latin{Graduale}, \latin{Ps.\ 33, 10 et 11} (Heb.\ 34). Missal p.~703, no.~3977.}

\begin{latinproper}
Timéte Dóminum, omnes Sancti eius: quóniam nihil deest timéntibus eum.
℣.~Inquiréntes autem Dóminum, non defícient omni bono.
\end{latinproper}

\begin{englishwitness}{Douay--Rheims (Challoner), Ps 33:10--11}
Fear the Lord, all ye his saints: for there is no want to them that fear him.
\dots\ but they that seek the Lord shall not be deprived of any good.
\end{englishwitness}

\noindent The Gradual's wording is not the Clementine Vulgate's. For
\latin{nihil deest}, ``nothing is lacking,'' the Clementine reads \latin{non
est inopia}, which the Douay translates ``there is no want''; for \latin{non
deficient}, ``shall not fail,'' it reads \latin{non minuentur}, ``shall not be
diminished,'' which the Douay translates ``shall not be deprived.'' The verse
leaves out the first half of verse~11, \latin{Divites eguerunt, et
esurierunt}, ``The rich have wanted, and have suffered hunger.'' St
Augustine's psalter reads \latin{nihil deest} at verse~10, as the Gradual does,
and \latin{non minuentur} at verse~11, as the Clementine does. The Vulgate and
the Douay count the psalm's title as verse~1, as the Hebrew does, so these
verses are 33:10--11 in both; English Bibles that leave the title unnumbered
call them 34:9--10.
\subsection{Alleluia}
\label{proper-alleluia}

\properlocus{Alleluia verse, \latin{Matth.\ 11, 28}. Missal p.~703, no.~3978.}

\begin{latinproper}
Allelúia, allelúia. ℣.~Veníte ad me, omnes qui laborátis et oneráti estis:
et ego refíciam vos. Allelúia.
\end{latinproper}

\begin{englishwitness}{Douay--Rheims (Challoner), Mt 11:28}
Come to me all you that labor and are burdened, and I will refresh you.
\end{englishwitness}

\noindent The verse is the whole of Mt 11:28 and agrees with the Clementine
text. It is the Lord's invitation after his thanks to the Father, who has hid
these things ``from the wise and prudent'' and revealed them ``to little ones''
(11:25); it goes on, ``Take up my yoke upon you, and learn of me, because I am
meek, and humble of heart'' (11:29). The feast never falls between Septuagesima
and Easter or in Paschaltide, so no Tract or paschal Alleluia is printed.

\subsection{Gospel}
\label{proper-gospel}

\properlocus{\latin{Sequentia sancti Evangelii secundum Matthaeum}, Mt 5:1--12a. Missal p.~703, no.~3979.}

\begin{latinproper}
In illo témpore: Videns Iesus turbas, ascéndit in montem, et cum sedísset,
accessérunt ad eum discípuli eius, et apériens os suum, docébat eos,
dicens: Beáti páuperes spíritu: quóniam ipsórum est regnum cælórum. Beáti
mites: quóniam ipsi possidébunt terram. Beáti qui lugent: quóniam ipsi
consolabúntur. Beáti qui esúriunt et sítiunt iustítiam: quóniam ipsi
saturabúntur. Beáti misericórdes: quóniam ipsi misericórdiam consequéntur.
Beáti mundo corde: quóniam ipsi Deum vidébunt. Beáti pacífici: quóniam
fílii Dei vocabúntur. Beáti qui persecutiónem patiúntur propter iustítiam:
quóniam ipsórum est regnum cælórum. Beáti estis cum maledíxerint vobis, et
persecúti vos fúerint, et díxerint omne malum advérsum vos, mentiéntes,
propter me: gaudéte et exsultáte, quóniam merces vestra copiósa est in
cælis.
\end{latinproper}

\begin{englishwitness}{Douay--Rheims (Challoner), Mt 5:1--12a}
And seeing the multitudes, he went up into a mountain, and when he was set
down, his disciples came unto him. And opening his mouth he taught them,
saying: Blessed are the poor in spirit: for theirs is the kingdom of heaven.
Blessed are the meek: for they shall possess the land. Blessed are they that
mourn: for they shall be comforted. Blessed are they that hunger and thirst
after justice: for they shall have their fill. Blessed are the merciful: for
they shall obtain mercy. Blessed are the clean of heart: for they shall see
God. Blessed are the peacemakers: for they shall be called the children of
God. Blessed are they that suffer persecution for justice' sake: for theirs is
the kingdom of heaven. Blessed are ye when they shall revile you, and
persecute you, and speak all that is evil against you, untruly, for my sake:
Be glad and rejoice for your reward is very great in heaven. \dots
\end{englishwitness}

\noindent The Missal's opening, \latin{In illo tempore: Videns Iesus turbas},
stands for the Vulgate's \latin{Videns autem Jesus turbas}, the Douay's ``And
seeing the multitudes''; the Lasance renders it ``At that time, Jesus seeing
the multitudes, went up into a mountain.'' The reading ends at the first half of
verse~12, \latin{in caelis}, and leaves out the second, ``For so they
persecuted the prophets that were before you.'' Every other word agrees with
the Clementine text. The Missal prints the
cue \latin{Credo} after the Gospel.

\subsection{Offertory}
\label{proper-offertory}

\properlocus{\latin{Antiphona ad Offertorium}, \latin{Sap.\ 3, 1, 2 et 3}. Missal p.~703, no.~3980.}

\begin{latinproper}
Iustórum ánimæ in manu Dei sunt, et non tanget illos torméntum malítiæ:
visi sunt óculis insipiéntium mori: illi autem sunt in pace, allelúia.
\end{latinproper}

\begin{englishwitness}{Douay--Rheims (Challoner), Wis 3:1--3}
\dots\ the souls of the just are in the hand of God, and the torment of death
shall not touch them. In the sight of the unwise they seemed to die: \dots\
but they are in peace.
\end{englishwitness}

\noindent The antiphon is a received liturgical form, not the verses as the
Vulgate prints them, and the Douay cannot render it word for word. Against
the Clementine it leaves out the opening \latin{autem}, ``But''; reads
\latin{tanget} for \latin{tangent}, a difference that does not show in
English; reads \latin{tormentum malitiae}, ``the torment of malice,'' for
\latin{tormentum mortis}, ``the torment of death''; leaves out ``and their
departure was taken for misery: And their going away from us, for utter
destruction'' (3:2b--3a); and adds \latin{alleluia}. The Lasance gives the
antiphon its own English: ``The souls of the just are in the hand of God, and
the torment of malice shall not touch them: in the sight of the unwise they
seemed to die, but they are in peace. Alleluia.''

The antiphon's wording has a history of its own. As a chant of martyrs'
Masses the verses were sung with \latin{malitiae} before the printed Missal:
so the Gregorian Antiphonary printed in Migne's \work{Patrologia Latina}
gives them, as a Communion and as a Gradual; so do the chant excerpts from
the Fonte Avellana codices; and so does the Sarum Missal in its common of many
martyrs, with the same omission as the antiphon, while the lesson from Wisdom
in the same common reads \latin{mortis}. The Venice printing of the Roman
Missal of 1570 reads \latin{tormentum mortis} at All Saints and has no
\latin{alleluia}, though the same book's Masses of martyrs read
\latin{malitiae}. The Roman printings of 1604, 1862 and 1889, and Fromage's
text of 1909, read \latin{malitiae} at the feast and have no \latin{alleluia}.
The earliest printing at hand that has the \latin{alleluia} is the Mame of
1922, conformed to the typical edition of 1920.
\subsection{Secret}
\label{proper-secret}

\properlocus{\latin{Secreta}. Missal p.~703, no.~3981.}

\begin{latinproper}
Múnera tibi, Dómine, nostræ devotiónis offérimus: quæ et pro cunctórum
tibi grata sint honóre iustórum, et nobis salutária, te miseránte,
reddántur. Per Dóminum.
\end{latinproper}

\begin{englishwitness}{Lasance, The New Roman Missal (1945), All Saints, Secret}
We offer Thee, O Lord, the gifts of our devotion; may they be made pleasing to
Thee for the honor of all the just and, by Thy mercy, be profitable to our
selves. Through our Lord.
\end{englishwitness}

\noindent The Lasance prints ``our selves'' as two words. The saints are named
in this prayer by the Offertory's word, \latin{iustorum}, ``the just.''

\subsection{Commemoration Secret}
\label{proper-commemoration-secret}

\properlocus{\latin{Secreta} of the Twenty-third Sunday after Pentecost. Missal p.~417, no.~1726.}

\begin{latinproper}
Pro nostræ servitútis augménto sacrifícium tibi, Dómine, laudis offérimus:
ut, quod imméritis contulísti, propítius exsequáris. Per Dóminum nostrum
Iesum Christum, Fílium tuum: Qui tecum vivit et regnat in unitáte.
\end{latinproper}

\begin{englishwitness}{Lasance, The New Roman Missal (1945), Twenty-third Sunday after Pentecost, Secret}
We offer Thee, O Lord, the sacrifice of praise for the increase of our
service, that Thou mayest graciously complete that which, for no merit of
ours, Thou hast granted unto us. Through our Lord.
\end{englishwitness}

\noindent It is said secretly after the feast's Secret, and its last words,
\latin{Per omnia saecula saeculorum}, are said aloud to begin the Preface
(RGMR 481). On the Sunday's own page the Missal directs a Preface after this
prayer; that direction belongs to the Mass of the Sunday, and a commemoration
does not bring a Preface of its own into the Mass in which it is made (RGMR
483).

\subsection{Preface}
\label{proper-preface}

\properlocus{No Preface is printed or directed with the feast. Either the Preface of the Most Holy Trinity, Missal p.~293, no.~1082, or the common Preface, p.~297, no.~1086.}

\noindent The feast has no Preface of its own, and its page names none, so
neither Preface's text is printed here. The general rubrics state that a Mass
takes its own Preface; that a Mass without one takes the Preface of the season
(\latin{de Tempore}); and that otherwise the common Preface is said (RGMR
482). The question is whether the Sunday's seasonal Preface reaches a feast of
the first class that has no Preface of its own and is said on a Sunday of the
second class.

One answer gives the Preface of the Most Holy Trinity. RGMR 494~b appoints
that Preface as the seasonal Preface on Sundays of the second class outside
Christmastide and Paschaltide, and the feast, lacking its own, would take the
seasonal one under RGMR 482. The Missal's direction at the Vigil of the
Nativity, which gives the common Preface but the Trinity Preface when the
Vigil falls on a Sunday (p.~16), is then an instance of the general rule. The
Vigil can fall only on the Fourth Sunday of Advent, which RGMR 494~b names in a
clause of its own; the Twenty-third Sunday after Pentecost falls under the
clause for the other Sundays of the second class, so the instance reaches All
Saints only if the two clauses work alike. The Ordo of the Institute of Christ
the King Sovereign Priest for 1 November 2026 directs the Trinity Preface;
that is one institute's direction, not a rubric of the Missal.

The other answer gives the common Preface. RGMR 494~b speaks of the Mass of
the Sunday and not of a feast said on a Sunday. The rubrics for Lent,
Passiontide, Easter, the Ascension and Pentecost extend their Prefaces in
express words to the other Masses of the season that lack a Preface of their
own (RGMR 486~b, 487~b, 489~b, 490~b, 493~b), and 494~b has no such words. The
Sunday is present only by its commemoration, and a commemoration brings no
Preface into the Mass (RGMR 483); the common Preface is said in Masses that
have no Preface of their own and are not to take the seasonal one (RGMR 498).
On this answer the Vigil's direction is an exception, printed because the
general rule did not reach the Vigil. The Missal's words admit both answers.
\subsection{Communion}
\label{proper-communion}

\properlocus{\latin{Antiphona ad Communionem}, \latin{Matth.\ 5, 8-10}. Missal p.~703, no.~3982.}

\begin{latinproper}
Beáti mundo corde, quóniam ipsi Deum vidébunt: beáti pacífici, quóniam
fílii Dei vocabúntur: beáti qui persecutiónem patiúntur propter
iustítiam, quóniam ipsórum est regnum cælórum.
\end{latinproper}

\begin{englishwitness}{Douay--Rheims (Challoner), Mt 5:8--10}
Blessed are the clean of heart: for they shall see God. Blessed are the
peacemakers: for they shall be called the children of God. Blessed are they
that suffer persecution for justice' sake: for theirs is the kingdom of
heaven.
\end{englishwitness}

\noindent The antiphon is Mt 5:8--10 word for word, the sixth, seventh and
eighth beatitudes of the Gospel just read.

\subsection{Postcommunion}
\label{proper-postcommunion}

\properlocus{\latin{Postcommunio}. Missal p.~703, no.~3983.}

\begin{latinproper}
Da, quǽsumus, Dómine, fidélibus pópulis ómnium Sanctórum semper
veneratióne lætári: et eórum perpétua supplicatióne muníri. Per Dóminum
nostrum.
\end{latinproper}

\begin{englishwitness}{Lasance, The New Roman Missal (1945), All Saints, Postcommunion}
Grant to Thy faithful people, we beseech Thee, O Lord, ever to rejoice in the
veneration of all Thy saints, and to be defended by their unceasing prayers.
Through our Lord.
\end{englishwitness}

\noindent The Latin's \latin{fidelibus populis} is plural, ``the faithful
peoples''; the Lasance gives the singular. The Mame printing ends the prayer
with the shorter cue \latin{Per Dominum}.

\subsection{Commemoration Postcommunion}
\label{proper-commemoration-postcommunion}

\properlocus{\latin{Postcommunio} of the Twenty-third Sunday after Pentecost. Missal p.~417, no.~1728.}

\begin{latinproper}
Quǽsumus, omnípotens Deus: ut, quos divína tríbuis participatióne
gaudére, humánis non sinas subiacére perículis. Per Dóminum nostrum Iesum
Christum, Fílium tuum: Qui tecum vivit et regnat in unitáte.
\end{latinproper}

\begin{englishwitness}{Lasance, The New Roman Missal (1945), Twenty-third Sunday after Pentecost, Postcommunion}
We pray Thee, O almighty God, that Thou wilt not suffer to be overcome by
human dangers those to whom Thou grantest participation in things that are
divine. Through our Lord.
\end{englishwitness}

\noindent The Lasance does not render the Latin's \latin{gaudere},
``to rejoice.'' Nothing else in the Mass is proper to the feast: it has no
Sequence, no proper \latin{Communicantes} or \latin{Hanc igitur}, and no
prayer over the people (RGMR 501).
